% Options for packages loaded elsewhere
\PassOptionsToPackage{unicode}{hyperref}
\PassOptionsToPackage{hyphens}{url}
\documentclass[
  11pt]{article}
\usepackage{xcolor}
\usepackage[margin=2.5cm]{geometry}
\usepackage{amsmath,amssymb}
\setcounter{secnumdepth}{-\maxdimen} % remove section numbering
\usepackage{iftex}
\ifPDFTeX
  \usepackage[T1]{fontenc}
  \usepackage[utf8]{inputenc}
  \usepackage{textcomp} % provide euro and other symbols
\else % if luatex or xetex
  \usepackage{unicode-math} % this also loads fontspec
  \defaultfontfeatures{Scale=MatchLowercase}
  \defaultfontfeatures[\rmfamily]{Ligatures=TeX,Scale=1}
\fi
\usepackage{lmodern}
\ifPDFTeX\else
  % xetex/luatex font selection
\fi
% Use upquote if available, for straight quotes in verbatim environments
\IfFileExists{upquote.sty}{\usepackage{upquote}}{}
\IfFileExists{microtype.sty}{% use microtype if available
  \usepackage[]{microtype}
  \UseMicrotypeSet[protrusion]{basicmath} % disable protrusion for tt fonts
}{}
\makeatletter
\@ifundefined{KOMAClassName}{% if non-KOMA class
  \IfFileExists{parskip.sty}{%
    \usepackage{parskip}
  }{% else
    \setlength{\parindent}{0pt}
    \setlength{\parskip}{6pt plus 2pt minus 1pt}}
}{% if KOMA class
  \KOMAoptions{parskip=half}}
\makeatother
\usepackage{longtable,booktabs,array}
\newcounter{none} % for unnumbered tables
\usepackage{calc} % for calculating minipage widths
% Correct order of tables after \paragraph or \subparagraph
\usepackage{etoolbox}
\makeatletter
\patchcmd\longtable{\par}{\if@noskipsec\mbox{}\fi\par}{}{}
\makeatother
% Allow footnotes in longtable head/foot
\IfFileExists{footnotehyper.sty}{\usepackage{footnotehyper}}{\usepackage{footnote}}
\makesavenoteenv{longtable}
\setlength{\emergencystretch}{3em} % prevent overfull lines
\providecommand{\tightlist}{%
  \setlength{\itemsep}{0pt}\setlength{\parskip}{0pt}}
\usepackage{bookmark}
\IfFileExists{xurl.sty}{\usepackage{xurl}}{} % add URL line breaks if available
\urlstyle{same}
\hypersetup{
  hidelinks,
  pdfcreator={LaTeX via pandoc}}

\author{}
\date{}

\begin{document}

\section{BH Thermodynamics Program Paper 9: Geometric Origin of the
0.036 Drift in the α Self-Consistency
Equation}\label{bh-thermodynamics-program-paper-9-geometric-origin-of-the-0.036-drift-in-the-ux3b1-self-consistency-equation}

\subsection{― Physical Interpretation via 2D Surface Vibration Modes in
a 4D
Hyperball}\label{physical-interpretation-via-2d-surface-vibration-modes-in-a-4d-hyperball}

\textbf{Author}: Noriaki Kihara \textbf{Affiliation}: WF System Co.,
Ltd. \textbf{ORCID}:
\href{https://orcid.org/0009-0004-6753-4020}{0009-0004-6753-4020}
\textbf{Version}: v1 \textbf{Date}: May 20, 2026 \textbf{License}: CC BY
4.0 \textbf{Concept DOI}:
\href{https://doi.org/10.5281/zenodo.20319436}{10.5281/zenodo.20319436}
\textbf{v1 Version DOI}:
\href{https://doi.org/10.5281/zenodo.20319437}{10.5281/zenodo.20319437}
\textbf{Zenodo page}: https://zenodo.org/records/20319437

\begin{center}\rule{0.5\linewidth}{0.5pt}\end{center}

\subsection{Nature of this Paper}\label{nature-of-this-paper}

\textbf{This is an observational/interpretive paper, not a
first-principles proof paper.}

We present one interpretation of the physical origin of the coefficient
\(\pi^2/2\) in the α identity \(\alpha^{-1} = 137 + (\pi^2/2)\alpha\) of
Paper 7 {[}BH7{]}, based on the uncertainty principle of standard
quantum mechanics and dimensional analysis of vibration modes.

This paper does \textbf{not} claim:

\begin{itemize}
\tightlist
\item
  That \(\pi^2/2\) admits no interpretation other than the one given
  here
\item
  That the high-precision residual of the observed value
  \(\alpha^{-1} = 137.035999...\) is fully derived by the present
  interpretation
\item
  To exclude other possible interpretations
\item
  To ``prove'', ``strengthen'', or ``extend'' the claims of Paper 7
\end{itemize}

What this paper offers is one way of reading. Paper 7's claims do not
depend on this paper. The numerical fact observed in Paper 7 (that
\(\alpha^{-1} = 137 + (\pi^2/2)\alpha\) holds at 8.7 ppb precision) is
unchanged by whether one adopts this interpretation.

\subsection{Scope}\label{scope}

The scope of this paper:

\begin{itemize}
\tightlist
\item
  \textbf{Subject}: \(\alpha^{-1}(Q^2 \to 0) = 137.036\) (geometric
  origin of α in the Thomson limit)
\item
  \textbf{Method}: Standard QM uncertainty principle + dimensional
  analysis of vibration modes in 4D space
\item
  \textbf{Conclusion}: The 0.036 drift can be interpreted as the
  zero-point expectation value of 2D surface vibration modes
\end{itemize}

This paper does \textbf{not} address:

\begin{itemize}
\tightlist
\item
  The high-energy running of α (\(\alpha^{-1}(M_Z) \approx 127.95\))
\item
  Geometric rederivation of vacuum polarization from Standard Model
  lepton/quark loops
\item
  Unification of coupling constants (α with \(\alpha_s\), \(\alpha_w\),
  etc.)
\item
  Connection with gravity
\end{itemize}

These are already well-described by Standard Model vacuum polarization
theory, and the present geometric interpretation takes the position of
providing a \textbf{boundary condition (value at \(Q^2 = 0\))} to that
framework. The geometric reformulation of the high-energy side is left
as future work (§8).

\subsection{Abstract}\label{abstract}

In Paper 7 {[}BH7{]}, the self-consistent equation using unit-cube
packing in a 4D ball,

\[\alpha^{-1} = 137 + \frac{\pi^2}{2}\alpha,\]

was reported to agree with the observed value
\(\alpha^{-1}(0) = 137.035999084\) at a relative precision of 8.7 ppb.
This paper presents the following physical interpretation of the
structure of this equation:

\begin{enumerate}
\def\labelenumi{\arabic{enumi}.}
\tightlist
\item
  \textbf{Although α is dimensionless, its physical contributions appear
  through area-dimensional quantities} (§2): Since the scattering cross
  section \(\sigma_T \propto \alpha^2\), the physical role of α is
  mediated by area.
\item
  \textbf{Classifying 4D space vibration modes by dimension, only 2D
  face modes can contribute to α under isotropic averaging} (§3):
  0D/1D/3D/4D modes carry directionality and vanish under averaging,
  while 2D area remains as a scalar quantity.
\item
  \textbf{The position phase space measure of 2D faces of 137 hypercubes
  distributed in a 4D unit ball equals \(\int_{B_4(1)} dV = \pi^2/2\)}
  (§4): While this coincides numerically with \(V_4(1)\), in our
  interpretation it is read as ``the integral of position degrees of
  freedom of 2D faces''.
\item
  \textbf{The W7 self-consistent equation
  \(\alpha^{-1} = 137 + (\pi^2/2)\alpha\) can be interpreted as a
  self-consistent correction from zero-point vibrations of 2D face modes
  of the 137 hypercubes} (§5).
\item
  \textbf{This structure is identical to the plaquette action of Wilson
  lattice gauge theory} (§6): The isomorphism shown in Paper 8 {[}BH8{]}
  is realized at the geometric level.
\item
  \textbf{There exists a theoretical lower bound on the observed value
  of α, arising from the width of the 2D mode distribution} (§7).
\end{enumerate}

Our interpretation is completed within the framework of standard quantum
mechanics and does not require new assumptions (discrete spacetime,
complex integer lattices, etc.).

\begin{center}\rule{0.5\linewidth}{0.5pt}\end{center}

\subsection{§1 Introduction}\label{introduction}

Paper 7 {[}BH7{]} presented the α self-consistent equation

\[\alpha^{-1} = N(1) + V_4(1)\alpha = 137 + \frac{\pi^2}{2}\alpha\]

using the 137 unit 4D cubes (tesseracts) inscribed in the 4D ball
\(B_4(R=3)\) and the volume of the 4D unit ball \(V_4(1) = \pi^2/2\).
This equation agrees with the observed value
\(\alpha^{-1}(0) = 137.035999084(21)\) at a precision of 8.7 ppb.

The structure of this self-consistent equation left the following
unresolved questions:

\textbf{Question 1}: Why does \(\pi^2/2\) appear as a coefficient of α?
- Numerically it coincides with \(V_4(1)\), but it is physically unclear
why the volume of a 4D unit ball should appear - Is it a mere
mathematical identity, or does it carry physical meaning?

\textbf{Question 2}: Geometric origin of the 0.036 drift - The geometric
structure giving the integer 137 (W6 paper) is clear - But the origin of
``0.036'' was only explained as a closure of the self-consistency

\textbf{Question 3}: Physical content of the isomorphism with Wilson
lattice gauge theory shown in W8 {[}BH8{]} - Structural correspondence
is established - But the physical content of the correspondence (why it
holds) remains unclear

This paper presents the following interpretation in response to these
questions:

\begin{quote}
\textbf{α is physically an area-dimensional quantity
(\(\sigma \propto \alpha^2\)), and only 2D face modes can contribute to
α among the 4D space vibration modes. The position phase space measure
of 2D face vibration modes of the 137 hypercubes gives \(\pi^2/2\),
which physically explains the coefficient structure of the W7
self-consistent equation.}
\end{quote}

This interpretation takes the standard QM uncertainty principle as its
starting point and requires no new assumptions.

\begin{center}\rule{0.5\linewidth}{0.5pt}\end{center}

\subsection{§2 Dimensional Analysis of α and Area
Dimensionality}\label{dimensional-analysis-of-ux3b1-and-area-dimensionality}

\subsubsection{2.1 α itself is
dimensionless}\label{ux3b1-itself-is-dimensionless}

The fine-structure constant α is defined in SI units as

\[\alpha = \frac{e^2}{4\pi\epsilon_0 \hbar c} \approx 7.297 \times 10^{-3}\]

and is dimensionless.

\subsubsection{2.2 Physical contributions appear through
α²}\label{physical-contributions-appear-through-ux3b1uxb2}

However, in many cases the quantities through which α is physically
involved appear as \(\alpha^2\):

\textbf{Thomson scattering cross section}:
\[\sigma_T = \frac{8\pi}{3} r_e^2 = \frac{8\pi}{3} \left(\alpha \cdot \frac{\hbar}{mc}\right)^2 \propto \alpha^2 \cdot [\text{length}]^2\]

\textbf{Rutherford scattering differential cross section}:
\[\frac{d\sigma}{d\Omega} \propto \alpha^2 \cdot [\text{length}]^2\]

\textbf{Bohr radius}:
\[a_0 = \frac{\hbar}{\alpha mc} \propto \alpha^{-1} \cdot [\text{length}]\]

\textbf{Classical electron radius}:
\[r_e = \alpha \cdot \frac{\hbar}{mc} \propto \alpha \cdot [\text{length}]\]

Thus, α is a quantity directly linked to \textbf{cross section (area
dimension)}.

\subsubsection{2.3 Correspondence with Geometric Structures in 4D
Space}\label{correspondence-with-geometric-structures-in-4d-space}

When considering geometric structures in 4D space to which α naturally
couples:

\begin{itemize}
\tightlist
\item
  σ (cross section) ∝ α² ∝ area
\item
  → α is a quantity that couples to \textbf{2-dimensional structures}
\end{itemize}

This is consistent with the fact that in Wilson lattice gauge theory,
the gauge coupling appears as the action coefficient of plaquettes (2D
faces, 2-forms).

\subsubsection{2.4 Starting Point of This
Paper}\label{starting-point-of-this-paper}

From the above:

\begin{quote}
\textbf{When considering contributions to α from vibration modes in 4D
space, the area dimensionality of α restricts the contributing modes to
those with ``2-dimensional area''.}
\end{quote}

This is the central observation of this paper.

\begin{center}\rule{0.5\linewidth}{0.5pt}\end{center}

\subsection{§3 Dimensional Classification of Vibration Modes in 4D
Space}\label{dimensional-classification-of-vibration-modes-in-4d-space}

\subsubsection{3.1 Vibration Modes of 137 Hypercubes and the
Circumscribed
4-Sphere}\label{vibration-modes-of-137-hypercubes-and-the-circumscribed-4-sphere}

The system treated in Paper 7: - 4D ball \(B_4(3)\) of radius \(R = 3\)
- 137 unit 4D cubes (tesseracts) packed inside

This system has the following vibration modes:

\textbf{Internal structure of one tesseract}:

{\def\LTcaptype{none} % do not increment counter
\begin{longtable}[]{@{}lll@{}}
\toprule\noalign{}
Structural element & Count & Independent vibration modes \\
\midrule\noalign{}
\endhead
\bottomrule\noalign{}
\endlastfoot
Vertex (0-cell) & 16 & 16 × 4 = 64 directions \\
Edge (1-cell) & 32 & Stretching 32, transverse, etc. \\
Face (2-cell, square) & 24 & Bulging, twisting \\
Solid (3-cell, cube) & 8 & Volume vibration \\
4-cell (whole tesseract) & 1 & Whole breathing mode \\
\end{longtable}
}

\textbf{Across all 137 tesseracts}: - Each structural element's
vibration multiplied by 137 - Interaction modes between adjacent cubes -
Phonon-like spectrum as a whole

\textbf{Circumscribed 4-sphere \(S^3\) (radius 3)}:

{\def\LTcaptype{none} % do not increment counter
\begin{longtable}[]{@{}
  >{\raggedright\arraybackslash}p{(\linewidth - 2\tabcolsep) * \real{0.5000}}
  >{\raggedright\arraybackslash}p{(\linewidth - 2\tabcolsep) * \real{0.5000}}@{}}
\toprule\noalign{}
\begin{minipage}[b]{\linewidth}\raggedright
Mode
\end{minipage} & \begin{minipage}[b]{\linewidth}\raggedright
Description
\end{minipage} \\
\midrule\noalign{}
\endhead
\bottomrule\noalign{}
\endlastfoot
Radial breathing & Vibration of R = 3 \\
Spherical harmonics \(Y_{\ell m n}\) on \(S^3\) & Angular vibrations
(\(\ell = 0, 1, 2, \ldots\)) \\
4-axis distortion & Sphere → ellipsoid deformation \\
\end{longtable}
}

\subsubsection{3.2 Survival Conditions under Isotropic
Averaging}\label{survival-conditions-under-isotropic-averaging}

Among these modes, what matters for contributions to α is the
\textbf{survival condition under isotropic averaging}.

α is dimensionless and scalar, invariant under SO(4) symmetry. Therefore
contributions to α are limited to \textbf{modes that do not vanish under
isotropic averaging (averaging over 4-axis directions)}.

{\def\LTcaptype{none} % do not increment counter
\begin{longtable}[]{@{}
  >{\raggedright\arraybackslash}p{(\linewidth - 4\tabcolsep) * \real{0.3333}}
  >{\raggedright\arraybackslash}p{(\linewidth - 4\tabcolsep) * \real{0.3333}}
  >{\raggedright\arraybackslash}p{(\linewidth - 4\tabcolsep) * \real{0.3333}}@{}}
\toprule\noalign{}
\begin{minipage}[b]{\linewidth}\raggedright
Mode dimension
\end{minipage} & \begin{minipage}[b]{\linewidth}\raggedright
Mathematical property
\end{minipage} & \begin{minipage}[b]{\linewidth}\raggedright
Result of isotropic averaging
\end{minipage} \\
\midrule\noalign{}
\endhead
\bottomrule\noalign{}
\endlastfoot
0D (point) & Vector position fluctuation & \textbf{Zero} (translation
symmetry) \\
1D (line segment) & Has direction & \textbf{Zero} (line symmetry, sign
flips on reversal) \\
\textbf{2D (face)} & \textbf{Area (scalar quantity, no direction)} &
\textbf{Survives} \\
3D (volume, oriented) & Volume element \(dV = dx \wedge dy \wedge dz\)
is pseudoscalar (has orientation) & \textbf{Zero} (sign flips on
reversal) \\
4D (4-volume) & 4-volume element is pseudoscalar & \textbf{Zero} \\
\end{longtable}
}

\textbf{Key fact}: \textbf{Area is a scalar quantity}, and area is
invariant under flipping the orientation of the face. In contrast,
length and volume (oriented volume elements) change sign under reversal
of orientation.

\subsubsection{3.3 Selection Principle}\label{selection-principle}

From the above:

\begin{quote}
\textbf{Among vibration modes in 4D space, only 2D face modes can give
significant contributions to α under isotropic averaging.}
\end{quote}

This is consistent with the area dimensionality of α (§2).

\subsubsection{3.4 Consistency with Wilson Lattice Gauge
Theory}\label{consistency-with-wilson-lattice-gauge-theory}

This conclusion agrees with the structure of \textbf{2-form gauge fields
(plaquette action)} in Wilson lattice gauge theory:

\begin{itemize}
\tightlist
\item
  0-form (scalar field): gauge-invariant point-like structure
\item
  1-form (vector potential \(A_\mu\)): gauge-dependent
\item
  \textbf{2-form (field strength \(F_{\mu\nu}\)): gauge-invariant,
  plaquette action \(S = \beta \sum_p \text{Re}\,\text{tr}(U_p)\)}
\item
  3-form, 4-form: higher-order structures
\end{itemize}

That α ``couples to faces'' is a standard fact in Wilson formalism.

\begin{center}\rule{0.5\linewidth}{0.5pt}\end{center}

\subsection{§4 Position Phase Space of 2D Face Modes of 137
Hypercubes}\label{position-phase-space-of-2d-face-modes-of-137-hypercubes}

\subsubsection{4.1 Position Degrees of Freedom of 2D Face
Modes}\label{position-degrees-of-freedom-of-2d-face-modes}

To fully specify a 2D face vibration mode requires:

\begin{itemize}
\tightlist
\item
  \textbf{Position}: Where the 2D face is located in the 4D ball
  \(B_4(R=3)\)
\item
  \textbf{Direction}: One of 6 possible 2-plane directions in 4D space
  (\(\binom{4}{2} = 6\))
\item
  \textbf{Amplitude}: Excitation level of the vibration
\end{itemize}

\subsubsection{4.2 Integration of Position Degrees of
Freedom}\label{integration-of-position-degrees-of-freedom}

For a fixed direction and amplitude, integrating the position degrees of
freedom of 2D faces within the 4D unit ball (scale-normalized) gives:

\[\text{2D face position phase space} = \int_{B_4(1)} dV = \frac{\pi^2}{2}\]

This coincides numerically with \textbf{\(V_4(1) = \pi^2/2\)}.

\subsubsection{4.3 Interpretation: Reading of Physical
Meaning}\label{interpretation-reading-of-physical-meaning}

{\def\LTcaptype{none} % do not increment counter
\begin{longtable}[]{@{}
  >{\raggedright\arraybackslash}p{(\linewidth - 2\tabcolsep) * \real{0.5000}}
  >{\raggedright\arraybackslash}p{(\linewidth - 2\tabcolsep) * \real{0.5000}}@{}}
\toprule\noalign{}
\begin{minipage}[b]{\linewidth}\raggedright
Interpretation
\end{minipage} & \begin{minipage}[b]{\linewidth}\raggedright
Physical meaning of \(\pi^2/2\)
\end{minipage} \\
\midrule\noalign{}
\endhead
\bottomrule\noalign{}
\endlastfoot
Old interpretation (W7) & Volume \(V_4(1)\) of 4D unit ball \\
\textbf{This paper} & \textbf{Measure of position degrees of freedom for
placing 2D face modes in 4D unit ball} \\
\end{longtable}
}

The two are \textbf{numerically identical, but physically different}.
Under our interpretation, \(\pi^2/2\) has a necessary reason to appear
as the coefficient of α: since α couples to 2D face modes, the volume of
their position phase space appears as the geometric coefficient of the
coupling constant.

\subsubsection{4.4 Direction and Amplitude
Contributions}\label{direction-and-amplitude-contributions}

The direction contribution from 6 possible 2-planes may appear as a
combinatorial factor separately. In this paper, we treat the position
integral as the dominant factor, and consider the directional factor as
an \(O(1)\) correction even if numerically effective.

The amplitude appears as α itself (excitation probability amplitude):

\[\text{Collective contribution of 2D face modes} = \underbrace{\frac{\pi^2}{2}}_{\text{position phase space}} \times \underbrace{\alpha}_{\text{amplitude per face}} = \frac{\pi^2}{2} \alpha\]

This physically explains the α coefficient term in the W7
self-consistent equation.

\begin{center}\rule{0.5\linewidth}{0.5pt}\end{center}

\subsection{§5 Physical Interpretation of the W7 Self-Consistent
Equation}\label{physical-interpretation-of-the-w7-self-consistent-equation}

\subsubsection{5.1 Structure of the
Equation}\label{structure-of-the-equation}

The self-consistent equation of Paper 7 {[}BH7{]}:

\[\alpha^{-1} = 137 + \frac{\pi^2}{2}\alpha\]

\subsubsection{5.2 Each Term under Our
Interpretation}\label{each-term-under-our-interpretation}

{\def\LTcaptype{none} % do not increment counter
\begin{longtable}[]{@{}
  >{\raggedright\arraybackslash}p{(\linewidth - 4\tabcolsep) * \real{0.3333}}
  >{\raggedright\arraybackslash}p{(\linewidth - 4\tabcolsep) * \real{0.3333}}
  >{\raggedright\arraybackslash}p{(\linewidth - 4\tabcolsep) * \real{0.3333}}@{}}
\toprule\noalign{}
\begin{minipage}[b]{\linewidth}\raggedright
Term
\end{minipage} & \begin{minipage}[b]{\linewidth}\raggedright
Mathematical content
\end{minipage} & \begin{minipage}[b]{\linewidth}\raggedright
Physical interpretation (this paper)
\end{minipage} \\
\midrule\noalign{}
\endhead
\bottomrule\noalign{}
\endlastfoot
\textbf{137} & Number of unit cubes packed in 4D ball \(B_4(3)\) &
\textbf{Contribution when all vibration modes are in the ground state
(lowest energy)}: integer-theoretic fact of 4D ℤ⁴ lattice \\
\textbf{\(\pi^2/2\)} & \(V_4(1)\), or in our interpretation, 2D face
mode position phase space & \textbf{Position degrees of freedom for
placing 2D face modes of 137 hypercubes in the 4D unit ball} \\
\textbf{α} (in coefficient term) & Coupling constant &
\textbf{Excitation amplitude of each 2D face mode} \\
\textbf{\((\pi^2/2)\alpha\)} & 0.036 (numerically) & \textbf{Collective
contribution of zero-point vibrations of 2D face modes} (from standard
QM uncertainty principle) \\
\end{longtable}
}

\subsubsection{5.3 Connection with Standard QM Uncertainty
Principle}\label{connection-with-standard-qm-uncertainty-principle}

Each 2D face mode has the following zero-point energy under the harmonic
oscillator approximation:

\[E_{\text{zero-point}}^{(i)} = \frac{1}{2}\hbar\omega_i\]

where \(\omega_i\) is the eigen-frequency of the i-th 2D face mode.

The collective contribution of all 2D face modes is:

\[\sum_i \frac{1}{2}\hbar\omega_i = (\text{position phase space measure}) \times (\text{amplitude density}) = \frac{\pi^2}{2} \cdot \alpha\]

Here the α on the right side represents the self-consistency condition
that the average excitation amplitude of each mode equals α itself.

\subsubsection{5.4 Self-Consistency}\label{self-consistency}

From the structure that α is the amplitude of 2D face modes and that the
zero-point contribution of 2D face modes provides a correction to
\(\alpha^{-1}\), the self-consistent equation

\[\alpha^{-1} = 137 + \frac{\pi^2}{2}\alpha \quad \Leftrightarrow \quad \frac{\pi^2}{2}\alpha^2 + 137\alpha - 1 = 0\]

is naturally derived.

\subsubsection{5.5 Numerical Agreement}\label{numerical-agreement}

The positive root of this equation:

\[\alpha^{-1} = 137.036010988\ldots\]

The relative error with the CODATA 2018 value
\(\alpha^{-1}(0) = 137.035999084(21)\) is \textbf{8.7 ppb}. This is the
fact already observed in Paper 7, and our interpretation provides a
\textbf{physical framework} explaining this numerical agreement.

\begin{center}\rule{0.5\linewidth}{0.5pt}\end{center}

\subsection{§6 Strengthening of the Isomorphism with Wilson Lattice
Gauge
Theory}\label{strengthening-of-the-isomorphism-with-wilson-lattice-gauge-theory}

Paper 8 {[}BH8{]} showed that there is a structural correspondence
between the W7 self-consistent equation and Wilson lattice gauge theory
based on Schläfli duality and \(B_4\) equivariance. Our 2D face mode
interpretation clarifies the \textbf{physical content} of this
correspondence.

\subsubsection{6.1 Correspondence}\label{correspondence}

{\def\LTcaptype{none} % do not increment counter
\begin{longtable}[]{@{}
  >{\raggedright\arraybackslash}p{(\linewidth - 2\tabcolsep) * \real{0.5000}}
  >{\raggedright\arraybackslash}p{(\linewidth - 2\tabcolsep) * \real{0.5000}}@{}}
\toprule\noalign{}
\begin{minipage}[b]{\linewidth}\raggedright
Wilson lattice gauge theory
\end{minipage} & \begin{minipage}[b]{\linewidth}\raggedright
W7 self-consistency (this paper's interpretation)
\end{minipage} \\
\midrule\noalign{}
\endhead
\bottomrule\noalign{}
\endlastfoot
Coupling constant \(\beta = 1/g^2\) & \(\alpha^{-1}\) \\
Plaquette (2-form structure) & 2D face vibration mode \\
Number of plaquettes \(\sum_p\) & \(\pi^2/2\) (position phase space
measure of 2D faces) \\
Plaquette action \(\beta \sum_p \text{Re}\,\text{tr}(U_p)\) &
\((\pi^2/2)\alpha\) \\
Cell (4D hard core) & 137 (number of unit cube packing) \\
Total action & \(\alpha^{-1} = 137 + (\pi^2/2)\alpha\) \\
\end{longtable}
}

\subsubsection{6.2 Clarification of Physical
Content}\label{clarification-of-physical-content}

The isomorphism in Paper 8 was established as a \textbf{structural
correspondence}, but with this paper it is \textbf{physically
positioned} as:

\begin{quote}
\textbf{The W7 self-consistent equation is the realization of the
plaquette action structure of Wilson lattice gauge theory on the 4D ball
geometry.}
\end{quote}

\subsubsection{6.3 Implications}\label{implications}

With this: - The correspondence between W7 and W8 is elevated from
``mathematical isomorphism'' to ``physical isomorphism'' - α coupling to
``2D faces'' aligns with the natural interpretation of Wilson formalism
- A roadmap appears for \textbf{geometric reconstruction} of vacuum
polarization in standard QFT

\begin{center}\rule{0.5\linewidth}{0.5pt}\end{center}

\subsection{§7 Theoretical Lower Bound on Observed Width of
α}\label{theoretical-lower-bound-on-observed-width-of-ux3b1}

\subsubsection{7.1 Distribution Width of Zero-Point
Vibrations}\label{distribution-width-of-zero-point-vibrations}

Each 2D face mode's zero-point vibration has a finite distribution width
from the standard QM uncertainty principle:

\[\Delta n_i^{\text{zero}} = \frac{1}{2}\]

(zero-point fluctuation of the photon number operator)

\subsubsection{7.2 Propagation to α}\label{propagation-to-ux3b1}

The fluctuation in the collective contribution of all 2D face modes is:

\[(\Delta \alpha^{-1})^2 \approx \sum_i \left(\frac{\partial \alpha^{-1}}{\partial n_i}\right)^2 (\Delta n_i)^2\]

A concrete calculation is beyond the scope of this paper (requires
determination of mode eigen-frequencies), but structurally:

\begin{quote}
\textbf{There exists a theoretical lower bound on the observed value of
α, arising from the width of the 2D face mode distribution.}
\end{quote}

\subsubsection{7.3 Correspondence with CODATA
Measurement}\label{correspondence-with-codata-measurement}

\[\alpha^{-1}(0) = 137.035999084(21)\]

The CODATA uncertainty \((21) \approx 2.1 \times 10^{-8}\) is
experimental precision, and whether our theoretical lower bound is even
smaller requires further consideration of mode frequencies.

\subsubsection{7.4 Implications}\label{implications-1}

By our interpretation: - α is not a ``sharp point'' but the ``center
value of a distribution'' - The observed 8.7 ppb deviation (W7) leaves
room for higher-order mode contributions - Future precision measurements
may make our theoretical lower bound observable

\begin{center}\rule{0.5\linewidth}{0.5pt}\end{center}

\subsection{§8 Open Problems: Extension to High
Energies}\label{open-problems-extension-to-high-energies}

Our interpretation \textbf{addresses only the Thomson limit
\(\alpha^{-1}(Q^2 \to 0) = 137.036\)}.

\subsubsection{8.1 Running of α in Standard
QED}\label{running-of-ux3b1-in-standard-qed}

In the Standard Model, α runs with the energy scale \(Q^2\):

\begin{itemize}
\tightlist
\item
  \(\alpha^{-1}(0) = 137.036\) (Thomson limit)
\item
  \(\alpha^{-1}(M_Z^2) \approx 127.95\) (Z boson mass scale)
\item
  The difference \(\approx 9.08\) units is explained by vacuum
  polarization from lepton/quark loops
\end{itemize}

This running is \textbf{well-described} by precision measurements of
standard QED, and this paper does not compete with it.

\subsubsection{8.2 Position of This Paper}\label{position-of-this-paper}

Our geometric interpretation provides a physical origin for the
\textbf{boundary condition (value at \(Q^2 = 0\))} of α. Running is left
to standard QED.

The two are hierarchically complementary:

{\def\LTcaptype{none} % do not increment counter
\begin{longtable}[]{@{}
  >{\raggedright\arraybackslash}p{(\linewidth - 4\tabcolsep) * \real{0.3333}}
  >{\raggedright\arraybackslash}p{(\linewidth - 4\tabcolsep) * \real{0.3333}}
  >{\raggedright\arraybackslash}p{(\linewidth - 4\tabcolsep) * \real{0.3333}}@{}}
\toprule\noalign{}
\begin{minipage}[b]{\linewidth}\raggedright
Layer
\end{minipage} & \begin{minipage}[b]{\linewidth}\raggedright
Content
\end{minipage} & \begin{minipage}[b]{\linewidth}\raggedright
Responsibility
\end{minipage} \\
\midrule\noalign{}
\endhead
\bottomrule\noalign{}
\endlastfoot
1 & 137 (integer, geometric invariant) & Paper 6 {[}BH6{]}: integer
theory of packing \\
2 & \((\pi^2/2)\alpha \approx 0.036\) & \textbf{This paper: 2D face
modes} \\
3 & \(\Delta\alpha_{SM}(Q^2)\) & Standard QED vacuum polarization \\
\end{longtable}
}

\subsubsection{8.3 Future Tasks}\label{future-tasks}

The relationship between layer 3 (standard QED vacuum polarization) and
our geometric interpretation leaves the following open questions:

\begin{enumerate}
\def\labelenumi{\arabic{enumi}.}
\tightlist
\item
  \textbf{Additional modes activated at high energies}: 3D and 4D
  vibration modes may start contributing to α at high energies
\item
  \textbf{Geometric counterparts of SM loops}: Whether each of electron
  loop, muon loop, quark loop corresponds to a specific geometric mode
\item
  \textbf{Geometric derivation of \(\alpha^{-1}(M_Z) \approx 128\)}:
  Whether our framework can be extended to high energies
\item
  \textbf{α behavior at GUT scale}: Geometric prediction of coupling
  constants at grand unification scale
\end{enumerate}

The geometric reformulation of these is left as future research.
\textbf{The claim of this paper (geometric origin of the Thomson limit)
holds independently of the mechanism of running}.

\begin{center}\rule{0.5\linewidth}{0.5pt}\end{center}

\subsection{§9 Conclusion}\label{conclusion}

This paper presented the following as a physical interpretation of the α
self-consistent equation \(\alpha^{-1} = 137 + (\pi^2/2)\alpha\) in
Paper 7 {[}BH7{]}:

\begin{enumerate}
\def\labelenumi{\arabic{enumi}.}
\tightlist
\item
  \textbf{α is dimensionless but physically area-dimensional}
  (\(\sigma \propto \alpha^2\))
\item
  \textbf{Among 4D space vibration modes, only 2D face modes can
  contribute to α under isotropic averaging}
\item
  \textbf{Position phase space measure of 2D faces of 137 hypercubes
  distributed in the 4D unit ball = \(\pi^2/2\)}
\item
  \textbf{The W7 self-consistent equation can be interpreted as a
  self-consistent contribution of zero-point vibrations of 2D face
  modes}
\item
  \textbf{This structure is identical to Wilson lattice gauge theory
  plaquette action} (physical content of W8 isomorphism)
\item
  \textbf{The observed value of α has a theoretical lower bound arising
  from the 2D mode distribution width}
\end{enumerate}

Our interpretation is completed within the framework of the standard QM
uncertainty principle, and requires no new assumptions (discrete
spacetime, extra dimensions, etc.).

\textbf{The scope of this paper} is only the Thomson limit of α. The
high-energy running is left to standard QED, and its geometric
reformulation is left as future tasks (§8).

\subsubsection{Position of This Paper}\label{position-of-this-paper-1}

Relationship among Paper 7 {[}BH7{]}, Paper 8 {[}BH8{]}, and this paper:

{\def\LTcaptype{none} % do not increment counter
\begin{longtable}[]{@{}
  >{\raggedright\arraybackslash}p{(\linewidth - 4\tabcolsep) * \real{0.3333}}
  >{\raggedright\arraybackslash}p{(\linewidth - 4\tabcolsep) * \real{0.3333}}
  >{\raggedright\arraybackslash}p{(\linewidth - 4\tabcolsep) * \real{0.3333}}@{}}
\toprule\noalign{}
\begin{minipage}[b]{\linewidth}\raggedright
Paper
\end{minipage} & \begin{minipage}[b]{\linewidth}\raggedright
Content
\end{minipage} & \begin{minipage}[b]{\linewidth}\raggedright
Area
\end{minipage} \\
\midrule\noalign{}
\endhead
\bottomrule\noalign{}
\endlastfoot
BH7 (α identity) & Discovery of \(\alpha^{-1} = 137 + (\pi^2/2)\alpha\)
& Algebraic observation \\
BH8 (Wilson isomorphism) & Proof of structural correspondence &
Mathematical correspondence \\
\textbf{This paper (BH9)} & \textbf{Physical interpretation via 2D face
modes} & \textbf{Physical content} \\
\end{longtable}
}

With this trilogy, the geometric origin of α is positioned in three
layers: \textbf{observation → structure → physical content}.

\begin{center}\rule{0.5\linewidth}{0.5pt}\end{center}

\subsection{References}\label{references}

{[}BH6{]} Kihara, N. (2026). \emph{Synthesis and the 4+1 to 3+1
Reduction}. Zenodo. Concept DOI: 10.5281/zenodo.19837597.

{[}BH7{]} Kihara, N. (2026). \emph{A Geometric Identity for the
Fine-Structure Constant: From the 4D Unit Ball Volume and its
Cube-Packing Deficit}. Zenodo. Concept DOI: 10.5281/zenodo.19869266.

{[}BH8{]} Kihara, N. (2026). \emph{Chain Complex Structure on the 4D
Hypercubic Lattice: Structural Correspondence between Kihara
Cube-Packing and Wilson Lattice Gauge Theory}. Zenodo. Concept DOI:
10.5281/zenodo.19880467.

{[}BH7-Supp{]} Kihara, N. (2026). \emph{Paper 7 Supplement: Geometric
Observation on the Second-Order Correction Term of the α Identity}.
Zenodo. Concept DOI: 10.5281/zenodo.19933729.

{[}Wilson1974{]} Wilson, K. G. (1974). \emph{Confinement of quarks}.
Phys. Rev.~D \textbf{10}, 2445.

{[}CODATA2018{]} Tiesinga, E. et al.~(2021). \emph{CODATA Recommended
Values of the Fundamental Physical Constants: 2018}. Rev.~Mod. Phys.
\textbf{93}, 025010.

\begin{center}\rule{0.5\linewidth}{0.5pt}\end{center}

\subsection{Revision History}\label{revision-history}

\begin{itemize}
\tightlist
\item
  \textbf{v1 (May 20, 2026)}: Initial version
\end{itemize}

\end{document}
